% Source: Wald GR, physical PDF pp. 280--281
%         (printed pp. 267--268).
% Coverage: Chapter 10 problems and end material.
% Figures/tables/footnotes: no figures, tables, or footnotes.
% Status: source-reviewed and compile-clean.
% Uncertainties: source “Gauss-Codacci” in problem 4 corrected to the standard
%                “Gauss--Codazzi,” matching the correction in Section 10.2.

\chapterproblems

\begin{enumerate}[label=\arabic*.,leftmargin=*,itemsep=1em]
\item Show that the inequality~\eqref{eq:10.1.16} holds for any subset, \(A\), of
\(\mathbb R^n\) satisfying the uniform interior cone condition by the following
argument (Cantor 1973; Adams 1975): Let \(Q\) denote the solid closed cone in
\(\mathbb R^n\) of height \(H\) and solid angle, \(\Omega\), with vertex at the origin.
Let \(\psi:\mathbb R\to\mathbb R\) be a \(C^\infty\) function with \(\psi(r)=1\) for
\(r<H/3\) and \(\psi(r)=0\) for \(r>2H/3\).
  \begin{enumerate}[label=\alph*),leftmargin=2.2em,itemsep=.45em]
  \item For any \(C^\infty\) function \(f:Q\to\mathbb R\), show that for all integers
  \(k\geq1\), we have
  \[
    f(0)=\frac{(-1)^k}{(k-1)!}
      \int_0^{R(\theta_0)}r^{k-1}
      \frac{\dd^k}{\dd r^k}[\psi(r)f(r,\theta_0)]\,\dd r,
  \]
  where \(r\) is the usual spherical radial coordinate of \(\mathbb R^n\), \(\theta_0\)
  denotes any fixed angle inside the cone, and \(R(\theta_0)\) is the largest value of
  \(r\) in the cone at angle \(\theta_0\), i.e., the integral is taken over the portion
  of any ray through the origin lying within the cone.

  \item By integrating the result of (a) over all angles inside the cone, show that
  \[
    f(0)=C_1\int_Q r^{k-n}\frac{\dd^k}{\dd r^k}(\psi f),
  \]
  where \(C_1\) is a constant and the proper volume element of \(Q\) is understood
  in the integral.

  % SOURCE ERRATUM: printed ``Schwartz inequality''; the named result is the
  % Cauchy--Schwarz inequality.
  \item Using the Cauchy--Schwarz inequality, show that for \(k>n/2\) we have
  \(|f(0)|\leq C\lVert f\rVert_{Q,k}\) and consequently, that
  equation~\eqref{eq:10.1.16} holds in the region \(A\).
  \end{enumerate}

\item Let \((M,g_{ab})\) be a globally hyperbolic spacetime with spacelike Cauchy
surface \(\Sigma\). Consider Maxwell's equations~\eqref{eq:4.3.12} and
\eqref{eq:4.3.13} in \((M,g_{ab})\) and define \(E_a\) and \(B_a\) on \(\Sigma\) by
equations~\eqref{eq:4.2.21} and \eqref{eq:4.2.22}, with \(v^a\) taken to be the unit
normal, \(n^a\), to \(\Sigma\). Note that \(E_an^a=B_an^a=0\).
  \begin{enumerate}[label=\alph*),leftmargin=2.2em,itemsep=.45em]
  \item Show that Maxwell's equations imply that \(D_aE^a=4\pi\rho\) and
  \(D_aB^a=0\) on \(\Sigma\), where \(\rho=-j^an_a\) and \(D_a\) is the derivative
  operator on \(\Sigma\). Note that the first relation implies that Gauss's law holds
  on \(\Sigma\).

  \item Show that the source-free (\(j^a=0\)) Maxwell's equations have a well posed
  initial value formulation in the sense that given \(E^a\) and \(B^a\) on \(\Sigma\)
  subject to the constraints \(D_aE^a=D_aB^a=0\), there exists a unique solution,
  \(F_{ab}\), of Maxwell's equations throughout \(M\) with these initial data and
  furthermore that this solution has the appropriate continuity and domain of
  dependence properties. To avoid ``patching'' arguments, you may assume global
  existence of a vector potential \(A_a\).
  \end{enumerate}

\item Let \(n^a\) be a unit (i.e., \(n^an_a=-1\)) hypersurface orthogonal vector
field. Define \(h_{ab}\) by equation~\eqref{eq:10.2.10}. Show that
\(h_a{}^c\nabla_cn_b=\frac12\mathcal L_nh_{ab}\).

\item Derive the Gauss--Codazzi relation~\eqref{eq:10.2.24}. (Hint: Evaluate the
left-hand side of Eq.~\eqref{eq:10.2.24} by using formulas for \(K_{ab}\) and
\(D_a\).)

\item Show explicitly from equation~\eqref{eq:10.2.26} that the components
\(G_{\mu\nu}n^\nu\) contain no second time derivatives of the metric components.

\item Use ``function counting'' arguments like those given at the end of this chapter
to conclude that the electromagnetic field has ``two degrees of freedom for each
point of space.''
\end{enumerate}

\endgroup
